\begin{pastexam}[label={px:nsysu-111-dmalg-07}]{中山資工所資安組 111 學年度 第 10 題}
（設計、分析）

(Algorithm points) {[10\%]} Given a knapsack with maximum capacity $W$ and a set $S$ consisting of $n$
items, where each item $i$ has weight $w_i$ and benefit value $b_i$. Please describe an algorithm based
on dynamic programming for packing the knapsack to achieve a maximum total benefit value of the picked
items and analyze the time complexity of the algorithm. ($w_i$, $b_i$, and $W$ are integer values.)
\end{pastexam}

\begin{pastsoln}
這是 0/1 背包問題（每件物品拿或不拿），做法照資結書第~18.5~節。 以下假設 $w_i \ge 1$、$W \ge 0$；
重量為 $0$ 的物品價值非負就直接拿，可以先處理掉。

\textbf{子問題。} $B[i, c]$：只能從物品 $1, \ldots, i$ 裡挑、總重不超過 $c$ 時，能得到的最大價值
（$0 \le i \le n$、$0 \le c \le W$）。 答案是 $B[n, W]$。

\textbf{遞迴式。}
\[
B[i, c] =
\begin{cases}
0, & i = 0, \\
B[i-1, c], & i \ge 1 \text{ 且 } w_i > c, \\
\max\bigl\{B[i-1, c],\; B[i-1, c - w_i] + b_i\bigr\}, & i \ge 1 \text{ 且 } w_i \le c.
\end{cases}
\]
理由：$\{1..i\}$、容量 $c$ 的最佳解如果不拿物品 $i$，它就是 $\{1..i-1\}$、容量 $c$ 的最佳解；如果拿，
扣掉物品 $i$ 後剩下的部分必須是 $\{1..i-1\}$、容量 $c - w_i$ 的最佳解，否則換掉它能得到更大的總價值
（最佳子結構）。 兩種情形取大的。

\textbf{演算法。} 第 $i$ 列只用到第 $i-1$ 列，由 $i = 0$ 往下逐列填表：
\begin{pseudo}[unbreakable]{0/1 背包}
\Function{Knapsack}{$w, b, n, W$}
    \For{$c \gets 0$ \To $W$}
        \State $B[0, c] \gets 0$
    \EndFor
    \For{$i \gets 1$ \To $n$}
        \For{$c \gets 0$ \To $W$}
            \State $B[i, c] \gets B[i - 1, c]$
            \If{$w_i \le c$ \AND $B[i - 1, c - w_i] + b_i > B[i, c]$}
                \State $B[i, c] \gets B[i - 1, c - w_i] + b_i$
            \EndIf
        \EndFor
    \EndFor
    \State $c \gets W$ \Comment{還原挑了哪些物品}
    \For{$i \gets n$ \DownTo $1$}
        \If{$B[i, c] \ne B[i - 1, c]$}
            \State 輸出物品 $i$；$c \gets c - w_i$
        \EndIf
    \EndFor
    \State \Return $B[n, W]$
\EndFunction
\end{pseudo}
還原時，$B[i, c] \ne B[i-1, c]$ 表示不拿物品 $i$ 就達不到這個值，所以物品 $i$ 一定在解裡。 手算的
小例子見政大資科所 112 學年度第 2 題（第~\pageref{px:nccu-112-dsa-02}~頁）。

\textbf{時間複雜度。} 表有 $(n+1)(W+1)$ 格，每格 $O(1)$；還原 $O(n)$。 總時間 $\Theta(nW)$，
空間 $\Theta(nW)$；只要最佳值時可以只留一列（$c$ 由大往小更新），空間降成 $\Theta(W)$。

$\Theta(nW)$ 看起來是多項式，其實不是：$W$ 只要 $\lceil \lg (W+1) \rceil$ 個位元就能寫下，$nW$ 對
輸入的位元數是指數的，這叫擬多項式時間（資結書第~18.5.5~節）。 $W$ 很大（例如 $2^{40}$）時這個方法
就不實用；0/1 背包的決策版本是 NP 完備的，目前沒有對位元數是多項式的演算法。
\end{pastsoln}
